\paragraph{Upper bound and its allocation certificate.}
Put
\[
 \widetilde\tau=\frac1n\sum_{i=1}^n Z_i,
 \qquad
 Z_i=\sum_{a\in S_c}\frac{c_a\mathbf 1\{A_i=a\}}{q_a^\star}
 \bigl(Y_i^{\mathrm{obs}}-1/2\bigr).
\]
The contrast is nonzero, so \(L_c>0\).  For every \(a\in S_c\), the prescribed allocation satisfies
\(q_a^\star=|c_a|/L_c>0\); hence every displayed division is valid.  Independent assignment and
\(\sum_a c_a=0\) give, for the fixed response type of unit \(i\),
\[
 \mathbb E_{\mathcal D}Z_i
 =\sum_{a\in S_c}c_a(t_{i,a}-1/2)
 =\sum_{a\in\mathcal A_K}c_at_{i,a}=:\mu_i.
\]
Consequently \(\mathbb E_{\mathcal D}\widetilde\tau=\tau_c(z)\).  Since every potential outcome is binary,
\((t_{i,a}-1/2)^2=1/4\), and therefore
\[
 \mathbb E_{\mathcal D}Z_i^2
 =\frac14\sum_{a\in S_c}\frac{c_a^2}{q_a^\star}
 =\frac14\sum_{a\in S_c}L_c|c_a|
 =\frac{L_c^2}{4}=C_0(c).
\]
The variables \(Z_1,\ldots,Z_n\) are independent under the stipulated design.  Thus
\[
 \mathbb E_{\mathcal D}\{\widetilde\tau-\tau_c(z)\}^2
 =\frac1{n^2}\sum_{i=1}^n\{C_0(c)-\mu_i^2\}
 \leq \frac{C_0(c)}n.
\]
The total positive coefficient mass and the absolute total negative coefficient mass both equal
\(L_c/2\), because the contrast coefficients sum to zero.  Hence, for every binary response vector,
\[
 -\frac{L_c}{2}\leq\sum_a c_at_{i,a}\leq\frac{L_c}{2},
 \qquad
 \tau_c(z)\in[-L_c/2,L_c/2].
\]
Metric projection onto this closed interval cannot increase squared distance from any point in the
interval.  Applying this pointwise to \(\widetilde\tau\) proves
\[
 \max_{z\in\mathcal T_K^n}
 R_n(\mathcal D,\widehat\tau_{\mathrm{cHT}}^\star;z)
 \leq \frac{C_0(c)}n,
\]
and also verifies that the projected rule takes values in the frozen estimator action space.

For completeness, this allocation is certified within the independent-allocation upper-bound
calculation.  If \(q_a>0\) for every \(a\in S_c\), then Cauchy--Schwarz and
\(\sum_{a\in S_c}q_a\leq1\) imply
\[
 \frac14\sum_{a\in S_c}\frac{c_a^2}{q_a}
 \geq\frac14\left(\sum_{a\in S_c}|c_a|\right)^2=C_0(c).
\]
Equality holds precisely when no probability is assigned outside \(S_c\) and
\(q_a=|c_a|/L_c\) on \(S_c\).

\paragraph{Direct threshold-prior converse.}
We give a Bayes lower bound that is uniform over all assignment laws, including dependent laws, and
over all possibly biased estimators.  Let \(s=|S_c|\), let \(\epsilon_n=n^{-1/4}/4\), and define
\[
 w(u)=\frac{15}{16}(1-u^2)^2\mathbf 1\{|u|\leq1\},
 \qquad
 J=\int_{-1}^1\frac{\{w'(u)\}^2}{w(u)}\,du<\infty.
\]
For each \(a\in S_c\), independently draw \(P_a\) with density
\[
 p\longmapsto \epsilon_n^{-1}w\!\left(\frac{p-1/2}{\epsilon_n}\right),
\]
and put \(P_a=1/2\) for \(a\notin S_c\).  Independently draw
\(U_1,\ldots,U_n\) identically and uniformly on \([0,1]\), and set
\[
 Y_i(a)=\mathbf 1\{U_i\leq P_a\}.
\]
After \((P,U)\) is drawn, this construction is a fixed binary schedule.  It therefore defines a
legitimate prior on the finite schedule space \(\mathcal T_K^n\), and it is independent of every
admissible assignment law.

Define
\[
 \vartheta(P)=\sum_{a\in S_c}c_aP_a,
 \qquad
 g_P(u)=\sum_{a\in S_c}c_a\mathbf 1\{u\leq P_a\}.
\]
Conditional on \(P\),
\[
 \tau_c(z)=\frac1n\sum_{i=1}^ng_P(U_i),
 \qquad
 \mathbb E\{\tau_c(z)\mid P\}=\vartheta(P).
\]
The function \(g_P\) vanishes outside \([\min_aP_a,\max_aP_a]\), since
\(\sum_a c_a=0\).  This interval has length at most \(2\epsilon_n\), and every partial sum of the
zero-sum contrast has absolute value at most \(L_c/2\), so \(|g_P|\leq L_c/2\).  It follows that
\[
 e_n^2:=\mathbb E\{\tau_c(z)-\vartheta(P)\}^2
 =\frac1n\mathbb E\{\operatorname{Var}(g_P(U_1)\mid P)\}
 \leq\frac{L_c^2\epsilon_n}{2n}.
\]
In particular, \(e_n=O(n^{-5/8})=o(n^{-1/2})\).

Fix an arbitrary design \(\mathcal D\) and write
\(s_a=\mathbb E_{\mathcal D}r_a\).  Conditional on \((P,A)\), the observed outcomes are independent
and satisfy \(Y_i^{\mathrm{obs}}\sim\operatorname{Bernoulli}(P_{A_i})\).  Symmetry of \(w\) and a
Taylor expansion on the support \([1/2-\epsilon_n,1/2+\epsilon_n]\) give
\[
 \kappa_n=\mathbb E\!\left\{\frac1{P_a(1-P_a)}\right\}
 =4+O(\epsilon_n^2).
\]
For \(a\in S_c\), let \(S_a\) denote the joint prior--likelihood score with respect to \(P_a\).
The score covariance matrix is diagonal, with
\[
 I_a:=\mathbb E(S_a^2)=\kappa_ns_a+\frac{J}{\epsilon_n^2}>0.
\]
Indeed, conditional Bernoulli information is \(r_a/\{P_a(1-P_a)\}\), the rescaled prior information
is \(J/\epsilon_n^2\), likelihood and prior scores are orthogonal, and distinct-arm likelihood scores
have zero conditional covariance even when the assignments are dependent.  The prior term makes every
\(I_a\) strictly positive, so neither exposure positivity nor a rank condition is required.

Let \(T=T(A,Y^{\mathrm{obs}})\) be any estimator in the frozen action space.  Both \(w\) and \(w'\)
vanish at the endpoints of their support.  Integration by parts in \(P_a\), with no boundary term,
therefore yields
\[
 \mathbb E(TS_a)=0,
 \qquad
 \mathbb E\{\vartheta(P)S_a\}=-c_a,
 \qquad
 \mathbb E[\{T-\vartheta(P)\}S_a]=c_a.
\]
For every vector \(v=(v_a)_{a\in S_c}\), Cauchy--Schwarz gives
\[
 \mathbb E\{T-\vartheta(P)\}^2
 \geq\frac{(v^\top c)^2}{\sum_{a\in S_c}v_a^2I_a}.
\]
Choosing \(v_a=c_a/I_a\), and then applying Cauchy--Schwarz to
\((|c_a|/\sqrt{I_a})_{a\in S_c}\) and \((\sqrt{I_a})_{a\in S_c}\), gives
\[
 \mathbb E\{T-\vartheta(P)\}^2
 \geq\sum_{a\in S_c}\frac{c_a^2}{I_a}
 \geq\frac{L_c^2}{\sum_{a\in S_c}I_a}
 \geq\frac{L_c^2}{\kappa_n n+sJ/\epsilon_n^2}
 =\frac{C_0(c)}n+o(n^{-1}).
\]
Here \(\sum_{a\in S_c}s_a\leq n\), \(K\), hence \(s\), is fixed, and
\(\epsilon_n^{-2}=16n^{1/2}=o(n)\).  All remainder bounds are independent of the chosen design and
estimator.

The reverse triangle inequality in \(L^2\) now implies
\[
 \left[\mathbb E\{T-\tau_c(z)\}^2\right]^{1/2}
 \geq
 \frac{L_c}{\sqrt{\kappa_n n+sJ/\epsilon_n^2}}-e_n.
\]
The first term is asymptotic to \(L_c/(2\sqrt n)\), whereas \(e_n=o(n^{-1/2})\).  Squaring the
positive part proves
\[
 \mathbb E\{T-\tau_c(z)\}^2\geq\frac{C_0(c)}n+o(n^{-1})
\]
under the prior, uniformly over \(\mathcal D\) and \(T\).  The maximum over deterministic schedules
dominates the prior-average risk.  Thus every sequence of admissible designs and possibly biased
estimators has asymptotic worst-case risk at least \(C_0(c)/n+o(n^{-1})\).  Combining this converse
with the finite-sample upper bound proves
\[
 \lim_{n\to\infty}n\rho_n(c)=C_0(c)=\frac14
 \left(\sum_{a\in\mathcal A_K}|c_a|\right)^2.
\]

\paragraph{Sharper embedded two-arm consequence.}
The established embedded two-arm converse strengthens the preceding first-order lower bound at every
finite \(n\): if \(\rho_{n,2}=\rho_n((1,-1))\), then
\[
 \rho_n(c)\geq C_0(c)\rho_{n,2}.
\]
Using the established two-arm expansion, this entails
\[
 \rho_n(c)\geq \frac{C_0(c)}n-C_0(c)C_A n^{-4/3}+o(n^{-4/3}),
 \qquad
 0\leq d_n(c)\leq C_0(c)\{n^{-1}-\rho_{n,2}\}.
\]
This implication is universal in \(c\).  Equality in the embedded comparison is established when
\(|S_c|=2\); for \(|S_c|\geq3\), only the displayed inequality is used, so no unsupported multi-arm
second-order equality is asserted.  Finally, the exact response-type-game theorem identifies the
labeled minimax value used throughout with \(G_n(c)\), confirming that neither the direct converse nor
the embedded converse restricts the arbitrary-design, arbitrary-estimator action space.